\section{Discussion}
A common four-input model-development procedure connected nominal-visit-matched affected skin area and serum measurements to EASI throughout this Chinese treatment cohort. Its strongest performance was in follow-up disease activity, where most estimates were within two EASI points. Using the same calculation across visits also allowed changes to be expressed directly in EASI units. This makes the formula straightforward to examine and reproduce: the research equation takes BSA and three serum concentrations without a treatment label or previous score. Same-day clinical/serum matching was not verified, so its use as a contemporaneous clinical measurement requires a defined sampling window.

The equation combines visible disease extent with measured serum information. BSA is an established extent component of EASI, which explains why it supplied most predictive accuracy in the ablation.\citep{hanifin2022,home_easi} CCL17, CCL26 and IL-5 provide biological measurements, but including laboratory inputs does not make the hybrid score observer-independent. Their defensible contribution here is to place an EASI-aligned estimate alongside treatment-dependent molecular changes. CCL17 and CCL26 declined in selected treatment comparisons, while IL-5 change accompanied itch change in one group. These observations address inflammatory response; they do not prove that three assays improve EASI estimation. Skin-biomarker findings during abrocitinib similarly show that molecular response and clinical correlation need not coincide, although tissue expression and serum concentrations are different measurements.\citep{guttmanyassky2024skin} Conditional coefficient signs, including the negative IL-5 term, are not mechanistic weights. A separate blood-only inflammation scale or a claim of molecular necessity would require evidence beyond the present fitted EASI reference.

The expanded-input comparisons offer a practical trade-off. Current itch modestly improved the hybrid follow-up estimate, and adult DLQI added some information in a smaller subset, at the cost of additional subjective inputs. Their performance does not justify relabelling the result as a purely objective score. The appropriate choice depends on whether the intended output is a short clinical estimate or an integrated clinical--serum assessment; the current data support the latter as an investigational measurement, with molecular added accuracy still unestablished across visits.

Individual monitoring is a distinct test from agreement at a single visit. The formula reproduced much of the ordering of clinical improvement, yet changes beginning at baseline were less precise than later follow-up changes. Higher observed baseline severity was associated with larger errors. Severity mix, prior treatment and reference quality cannot be separated as causes by these data. Lesion intensity can vary despite similar affected area. Chinese serum endotypes and recent severity-linked proteomic findings also show broader inflammatory and epithelial-associated heterogeneity than a three-marker panel can represent.\citep{wu2023,lally2026} These observations motivate testing additional lesion information in severe disease, but they do not identify the cause of individual errors or establish that more blood markers will reduce them.

The p-EASI studies provide the closest precedent for a usable serum-to-EASI calculation.\citep{thijs2017,thijs2019,bakker2020} Their correlations and classification results encouraged molecular estimation; the present study adds explicit absolute errors, their distribution and within-person response checks in a local treatment cohort. Their use of serum alone differs from the present BSA--serum inputs. Measurement procedures can affect calibration, and both assay compatibility and case mix matter when moving a prediction rule between settings.\citep{luijken2020,riley2024validation} The source-table local-application analyses therefore provide neither a fair superiority test nor evidence that population ethnicity explains a performance difference.

Treatment context matters both for interpreting biomarkers and for comparing clinical outcomes. Patients receiving the higher initial abrocitinib dose had greater starting severity; subsequent reductions were selected after clinical response. Baseline adjustment and explicit maintenance pathways therefore address different questions. This distinction is consistent with randomized abrocitinib induction/withdrawal designs and observational dupilumab tapering studies, which define eligibility for maintenance or reduction before evaluating its success.\citep{blauvelt2022regimen,spekhorst2023} The present results describe adjusted initial-strategy associations and observed later management, rather than the effect of assigning each patient an alternative dose. Differential missingness, endpoint-specific residual imbalance and response-selected follow-up limit these comparisons even after baseline adjustment. They also preserve itch and quality-of-life outcomes as complementary clinical domains; neither is adequately recovered by the same small serum--BSA panel.

The proposed management concept keeps the research EASI estimate alongside measured skin findings, symptoms, quality of life and safety (Fig.~\ref{fig:Figure1}). It does not supply automatic escalation, reduction or stopping thresholds. Independent evaluation should test the frozen equation, assay procedure and a prespecified clinical decision together. Patient-level validation prevented a person's repeated visits from appearing in both training and testing, but the molecular panel had been selected earlier and week 24 had already been inspected in prior project analyses. The current evidence consequently supports within-cohort performance of the development procedure, with a clearly defined research formula for the next evaluation. Patient bootstrap intervals condition on saved predictions and omit model-development uncertainty; no independent validation of the single final formula has been performed. Follow-up overestimation and calibration slopes below one also argue against substituting estimates for clinical examination. Adequate numbers of active severe cases, independent clinical assessment and a prespecified tolerance for individual errors will determine whether the same simple calculation can support decisions across the full disease range.
